% ============================================================
%  Interpreting FCFS Appointment Simulation Results
%  Under the Current Formulation
%  Metric behavior, access losses, and model checks
% ============================================================

\documentclass[11pt]{article}

\usepackage[margin=0.9in]{geometry}
\usepackage{amsmath}
\usepackage{amssymb}
\usepackage{booktabs}
\usepackage{graphicx}
\usepackage{float}
\usepackage[hidelinks]{hyperref}
\usepackage{array}
\usepackage[table]{xcolor}
\usepackage{microtype}
\usepackage{parskip}
\usepackage{enumitem}
\usepackage{mdframed}
\usepackage{tabularx}

\graphicspath{%
  {../../figures/}%
  {../../../../../outputs/class1_balking/figures/}%
  {../../../../../outputs/class1_balking_threshold/figures/}%
}

\sloppy
\setcounter{tocdepth}{2}

\newcommand{\mainfigure}[3]{%
\begin{figure}[H]
    \centering
    \includegraphics[width=\textwidth]{#1}
    \caption{#2}
    \label{#3}
\end{figure}%
}

\newmdenv[backgroundcolor=yellow!12, linecolor=gray!60, linewidth=0.6pt,
          innertopmargin=6pt, innerbottommargin=6pt]{checkenv}

\title{Interpreting FCFS Appointment Simulation Results\\
Under the Current Formulation\\[0.4em]
\large Metric behavior, access losses, and model checks}
\author{CUIMC Appointment Simulation}
\date{May 2026}

\begin{document}
\maketitle

\begin{abstract}
\noindent
This note interprets simulation outputs from a discrete-event FCFS appointment
model. The goal is to understand what the outputs show and whether they are
mechanically plausible under the current formulation---not to make empirical
claims about real clinics.

The central finding is that utilization and served rate separate sharply.
At baseline ($\lambda = 100$ arrivals per day, 32 slots per day, 14-day
horizon), average utilization is 0.839 while overall served rate is 0.269.
This separation is mechanically plausible because the two metrics have
different denominators: utilization is normalized by slots, served rate by
arrivals. The baseline is structurally oversubscribed. Patients are indeed served
on future days (the horizon extends to $H = 14$ days), but each service day
has exactly 32 slots regardless of which arrival cohort fills them. Over
any measurement window of $D$ days, at most $32D$ visits can be completed
against approximately $100D$ arrivals, so served rate cannot exceed
$32/100 = 0.32$.

Arrival outcome decomposition at baseline shows that balking (35.6\% of
arrivals) and cancellation (32.3\%) are the dominant loss channels, not
horizon saturation: no-offer is~0 at baseline demand. The high cancellation
fraction is mechanically consistent with the long average booking delay
(8.35~days) combined with a 10\%~daily cancellation probability.

Mean offered wait is preferable to mean accepted wait as a delay measure,
because offered delay is recorded before the balking decision and therefore
includes patients who ultimately reject the offer. Offered wait still excludes
no-offer patients, so it must be read together with the no-offer share---which
is zero at baseline but grows above $\lambda \approx 110$.
\end{abstract}

\tableofcontents
\newpage

%% ============================================================
\section{Simulation Mechanics Relevant for Interpretation}
\label{sec:mechanics}
%% ============================================================

The following mechanics directly affect how outputs should be read. The full
formulation is in the simulation documentation.

\paragraph{Finite booking horizon and FCFS pooling.}
Each day, arrivals for both classes are pooled and randomly permuted into one
FCFS sequence. Each patient is offered the earliest available slot within the
$H$-day horizon (including same-day, $r=0$). Because FCFS operates on a
permuted joint sequence, class labels are irrelevant to slot assignment:
neither class is prioritized. Class disparities in outcomes arise only from
differences in arrival rates or behavioral parameters.

\paragraph{Offered delay $\tau$ recorded before balking.}
When a valid slot exists, the offered delay $\tau$ is recorded before the
patient decides whether to accept. Mean offered wait therefore includes
patients who ultimately balk. Mean accepted wait excludes them, which biases
it downward when high-$\tau$ patients disproportionately balk.

\paragraph{No-show based on original $\tau$.}
No-show probability is a function of the original offered delay $\tau$, not of
the remaining residual day $r$ at service time. A patient who accepted a slot
9~days out carries their high-delay no-show probability regardless of when they
end up being seen.

\paragraph{Future cancellations only; no rebooking of no-show slots.}
Cancellations apply only to future appointments ($r \geq 1$). Same-day
cancellations are not modeled. Slots vacated by no-shows are not
rebooked---once a patient no-shows, that slot is lost.

\paragraph{No-offer rule.}
If the horizon is full when a patient reaches the front of the FCFS queue,
that patient and all remaining arrivals for that day are counted as
\texttt{no\_offer}. This means a single saturation event affects the entire
remainder of the day's queue, which could introduce daily-order artifacts.

\paragraph{Utilization denominator.}
Utilization counts completed visits per available slot, not scheduled
appointments per slot. No-shows, cancellations, and unbooked slots all reduce
utilization. A high-utilization day is one where most slots became completed
visits.

%% ============================================================
\section{Baseline Setup and Diagnosis}
\label{sec:baseline}
%% ============================================================

\subsection{Parameter Setup}

The baseline (\texttt{configs/baseline.yaml}) parameters are listed in
Figure~\ref{fig:baseline-params-and-step-functions}. The key feature is the demand-to-capacity ratio:
100~arrivals per day against 32~slots per day is a 3:1 oversubscription,
well above any realistic clinic load. The purpose is to probe how the
metrics behave across a wide range, not to model a plausible clinic.

\begin{figure}[H]
\centering

\begin{minipage}[t]{0.4\textwidth}
\vspace{0pt}
\renewcommand{\arraystretch}{1.15}
\small
\begin{tabularx}{\linewidth}{Xr}
\toprule
Parameter & Value \\
\midrule
Class 1 arrival rate $\lambda_1$ (per day) & 50 \\
Class 2 arrival rate $\lambda_2$ (per day) & 50 \\
Total arrival rate $\lambda$ (per day) & 100 \\
Slots per day $S$ & 32 \\
Booking horizon $H$ (days) & 14 \\
Cancellation probability $\phi$ & 0.10 \\
Balking threshold $\theta^{\text{balk}}$ (days) & 9 \\
High-delay balking prob.\ $b^{\text{high}}$ & 0.50 \\
Low-delay balking prob.\ $b^{\text{low}}$ & 0.00 \\
No-show threshold $\theta^{\text{no-show}}$ (days) & 6 \\
High-delay no-show prob.\ $\xi^{\text{high}}$ & 0.30 \\
Low-delay no-show prob.\ $\xi^{\text{low}}$ & 0.00 \\
\bottomrule
\end{tabularx}
\end{minipage}%
\hfill
\begin{minipage}[t]{0.58\textwidth}
\vspace{0pt}
\centering
\includegraphics[width=\linewidth,height=0.82\linewidth,keepaspectratio]{baseline_step_functions.png}
\end{minipage}

\caption{Left: baseline parameter values (\texttt{configs/baseline.yaml}),
symmetric across both classes; measurement window 365 days after a 30-day
burn-in. Right: balking and no-show step functions at baseline. Both
probabilities are~0 below their respective thresholds and jump to 0.50
(balking) and 0.30 (no-show) above them.}
\label{fig:baseline-params-and-step-functions}
\end{figure}

\subsection{Aggregate Metrics at Baseline}

Metrics are averaged over a single run of 365 measurement days (seed 11),
preceded by a 30-day burn-in to reach steady state and followed by a
10-day cooldown to allow late bookings to resolve. All figures in
Table~\ref{tab:baseline} are point estimates from that single run; the
sensitivity appendix uses multiple seeds (2027--2029 for slices, 4101--4102
for the regression screen) to assess stability and the influence of each parameter.

\begin{table}[H]
\centering
\renewcommand{\arraystretch}{1.15}
\begin{tabular}{lr}
\toprule
Metric & Value \\
\midrule
Average utilization & 0.839 \\
Overall served rate & 0.269 \\
Mean offered wait (days) & 9.30 \\
Mean accepted wait (days) & 8.35 \\
Overall balking rate & 0.356 \\
Class served-rate gap ($\Delta_{\text{access}}$) & 0.001 \\
\bottomrule
\end{tabular}
\caption{Baseline aggregate metrics. The near-zero class gap is expected:
both classes have identical behavioral parameters and symmetric arrival rates,
so class labels are exchangeable under pooled FCFS.}
\label{tab:baseline}
\end{table}

\subsection{The Utilization--Access Gap}

\[
  \text{average utilization}
  = \frac{\text{completed visits}}{\text{available slots}},
  \qquad
  \text{overall served rate}
  = \frac{\text{completed visits}}{\text{arrivals}}.
\]

Utilization~(0.839) and served rate~(0.269) differ by a factor of roughly
3.1. This is mechanically expected because the denominators differ.
Utilization is normalized by supply; served rate is normalized by demand.
When demand far exceeds supply, served rate is bounded by the
supply-to-demand ratio; utilization is not.

At baseline, $100/32 = 3.125$~arrivals per slot. The ceiling argument is
a flow constraint: patients are served on future days (via the booking
horizon), but each service day has exactly 32 slots regardless of which
arrival cohort fills them. Over $D$ days at most $32D$ visits can complete
against $100D$ arrivals, so served rate cannot exceed $1/3.125 = 0.32$.
The observed served rate of 0.269 falls below this ceiling because
no-shows consume slots without producing completed visits.

The simulation implies, under these assumptions, that a facility-level
utilization measure would give a misleading picture of patient access at
this demand-to-capacity ratio.

%% ============================================================
\section{Arrival Outcome Decomposition}
\label{sec:decomp}
%% ============================================================

\subsection{Loss Channels at Baseline}

Every arrival ends up in exactly one outcome state. At baseline:

\begin{table}[H]
\centering
\renewcommand{\arraystretch}{1.12}
\begin{tabular}{lrp{0.52\textwidth}}
\toprule
Outcome & Share of arrivals & Interpretation \\
\midrule
Served & 0.269 & Completed visit \\
Balked & 0.356 & Received an offer, rejected the delay \\
Canceled & 0.323 & Booked, but canceled before service day \\
No-show & 0.052 & Kept booking, did not attend \\
No-offer & 0.000 & Horizon full; no slot offered \\
Unresolved & $<$0.001 & Booked but not resolved before simulation end \\
\midrule
Total & 1.000 & \\
\bottomrule
\end{tabular}
\caption{Arrival outcome decomposition at baseline ($\lambda = 100$).
No-offer is zero: the horizon does not completely fill at baseline demand
because cancellations continuously free future capacity.}
\label{tab:decomp}
\end{table}

Two findings stand out. First, the no-offer rate is zero even at 100
arrivals/day. The horizon does not saturate because cancellations
(10\%/day for each future appointment) continuously free slots. However,
cancellation frees far-future capacity, which is then offered to new
arrivals at long delays---reinforcing the same long-delay dynamic that
drives balking.

Second, the 32.3\% cancellation rate (as a share of arrivals) is not
self-evidently wrong. Of the 64.4\% of arrivals that book an appointment,
roughly half ultimately cancel. With an average booking delay of 8.35~days
and a 10\%~daily cancellation probability, the implied cumulative
cancellation probability is $1 - 0.9^{8.35} \approx 0.59$. The observed
cancellation rate among booked patients is in the same range, though
somewhat lower---possibly because some booked patients face fewer
cancellation-check days when scheduled soon.

Inspection of \texttt{simulation/engine.py} confirms that the engine applies a
fresh Bernoulli draw to every future appointment on each simulated day, so the
cumulative argument is consistent with the implementation.

\subsection{How Loss Channels Shift with Demand}

Figure~\ref{fig:decomp-full} shows the arrival outcome decomposition as total
daily arrivals vary from 30 to 170. At low demand ($\lambda \lesssim 50$),
virtually all arrivals are served. As demand rises above the 32-slot daily
capacity, the served rate falls. The no-offer rate remains zero up to
$\lambda \approx 110$, at which point the horizon begins to saturate.

The transition from a cancellation- and balking-dominated regime (below
$\lambda \approx 110$) to a horizon-saturation regime (above) is a
structural feature of the model. It suggests that reporting the no-offer
share alongside served rate and offered wait is necessary to understand
which failure mode dominates in a given parameter setting.

\mainfigure{fcfs_arrival_outcome_decomposition.png}{%
Arrival outcome decomposition as total daily arrivals vary from 30 to 170.
At baseline ($\lambda = 100$, marked by a dashed vertical line), balking and
cancellation are the dominant loss channels; no-offer is zero. No-offer
becomes substantial only above $\lambda \approx 110$.}{fig:decomp-full}

%% ============================================================
\section{Cross-Metric Diagnostics}
\label{sec:crossmetric}
%% ============================================================

Figure~\ref{fig:main-panel} shows three diagnostics that together
capture the core structure of the simulation output.

\mainfigure{research_note_main_panel.png}{%
\textit{Left:} Arrival outcome decomposition as total daily arrivals
vary. At baseline ($\lambda = 100$, dashed line), balking and
cancellation are the dominant loss channels; no-offer is
zero. No-offer appears only above $\lambda \approx 110$.
\textit{Center:} Utilization--access frontier. The plot reveals two
regimes separated by the onset of balking. When balking is active
($\lambda \geq 60$, the main cluster), utilization is pinned near~0.84
while served rate falls from~0.45 to~0.16. When balking is absent
($\lambda \leq 50$, the three isolated points), both metrics co-move.
\textit{Right:} Offered vs.\ accepted wait. Points below the~45$^\circ$
line indicate that accepted wait understates the delay patients face.
The gap widens with balking rate (red intensity).}{fig:main-panel}

\textbf{1.~Utilization--access frontier (center panel).}
The center panel shows two distinct regimes, cleanly separated by whether
balking is active. In the balking regime ($\lambda \geq 60$), utilization is
pinned near~0.84 while served rate collapses from~0.45 to~0.16: as one
patient balks or cancels, a waiting patient fills the slot, sustaining the
slot-fill rate even as fewer arrivals ultimately receive service. In the
no-balking regime ($\lambda \leq 50$), that queue does not exist, and the
two metrics co-move. The apparent decoupling in the cluster is therefore
not a general property of the metrics; it is a consequence of balking
providing a surplus of waiters to absorb every freed slot.

\textbf{2.~Offered vs.\ accepted wait (right panel).}
All points fall below the 45$^\circ$ line: accepted wait consistently
understates offered wait. The gap widens as balking rate increases (red
intensity). This is mechanically consistent: patients who balk tend to have
above-average $\tau$, so removing them from the accepted sample pulls the
mean down. Any policy or parameter change that raises balking will make
accepted wait appear to improve while offered wait and served rate both
worsen.

\textbf{3.~No-offer and the boundary of offered-wait validity.}
At baseline, no-offer is zero, so offered wait covers all patients who
received any delay information. Above $\lambda \approx 110$, a growing
share of arrivals receive no offer and are excluded from offered wait
entirely. At $\lambda = 170$, the no-offer share reaches approximately
34.6\% of arrivals. In that regime, offered wait alone misses the
worst-off patients.

%% ============================================================
\section{Metric Hierarchy Under This Formulation}
\label{sec:metrics}
%% ============================================================

\begin{table}[H]
\centering
\renewcommand{\arraystretch}{1.2}
\begin{tabular}{p{0.20\textwidth}p{0.22\textwidth}p{0.24\textwidth}p{0.24\textwidth}}
\toprule
Metric & What it measures & Mechanical interpretation & Main caveat \\
\midrule
Overall served rate & Arrivals $\to$ completed visits & Primary access measure; bounded above by $S/\lambda$ & Does not decompose loss channels \\
Average utilization & Completed visits per available slot & Capacity-use measure; decoupled from served rate only when balking is active and a surplus of waiters refills every freed slot & High value does not imply good access \\
Mean offered wait & Delay among patients who received any offer & Less selective than accepted wait; includes patients who then balk & Excludes no-offer patients entirely \\
Mean accepted wait & Delay among patients who accepted bookings & Conditional; pulled down when high-delay patients balk & Can fall as access worsens \\
Balking rate & Share of offered patients who reject & Congestion diagnostic & Not a success metric \\
Class served-rate gap & Class 1 served rate minus Class 2 & Meaningful only when classes differ in behavior or arrivals & Near-zero in symmetric baseline \\
\bottomrule
\end{tabular}
\caption{Metric interpretation under the current FCFS formulation.}
\label{tab:metrics}
\end{table}

\subsection{Balking-Mediated Decoupling of Utilization and Access}

The center panel of Figure~\ref{fig:main-panel} separates into two regimes
with a clean boundary at the onset of balking. Every point with $\lambda \geq
60$ has a positive balking rate (ranging from~4.6\% to~36.1\%) and
utilization constrained to~0.839--0.841; every point with $\lambda \leq 50$
has zero balking, and both utilization and served rate vary freely.

In the balking regime, utilization is pinned near~0.84 while served rate falls
monotonically from~0.45 to~0.16 as demand rises from~60 to~170. This is a
direct consequence of how cancelled slots are filled: as long as a surplus of
balking-tolerant patients is waiting, each freed slot is immediately claimed,
keeping the slot-fill rate near a no-show-adjusted ceiling regardless of how
many arrivals ultimately fail to receive service. The apparent decoupling of
utilization from served rate holds only within this regime---it is a property
of the balking mechanism, not of the metrics themselves.

The three no-balking points ($\lambda \in \{30, 40, 50\}$) tell a different
story. At $\lambda = 30$, demand is below daily capacity; virtually all
arrivals are served (served rate~$= 0.998$) and utilization reaches~$0.943$,
though the horizon does not fully saturate from a single day's arrivals.
At $\lambda = 40$, slots are completely absorbed (utilization~$= 1.000$), but
20\% of arrivals cancel before their service day; with no queue to refill
those slots, served rate falls to~$0.800$. At $\lambda = 50$, demand exceeds
daily capacity yet no patient balks; a growing cancellation chain that goes
unrefilled pulls utilization back to~$0.903$ and served rate to~$0.576$.
Across all three, utilization and served rate co-move rather than decouple,
confirming that balking---not a structural property of the model---is what
sustains high slot fill under declining access.

Figure~\ref{fig:no-balk-delay} extends the comparison to $\lambda \leq 110$
and plots both mean offered and mean accepted delay against the two
behavioral thresholds. In the no-balking regime ($\lambda \leq 50$) the two
series coincide: every offer is accepted, so accepted and offered delay are
the same by construction. Mean accepted delay rises from $0.22$~days at
$\lambda = 30$ to $5.50$~days at $\lambda = 50$---in all three cases
strictly below the balking threshold ($\theta^{\text{balk}} = 9$~days),
confirming that zero balking is self-consistent. The $\lambda = 50$ point
sits just below the no-show threshold ($\theta^{\text{no-show}} = 6$~days),
however, which is why cancellation losses are elevated there: although the
mean does not breach~$6$~days, a non-trivial share of the delay
distribution does.

The two series split at exactly $\lambda = 60$, the onset of balking.
From that point onward offered delay continues to rise steeply while accepted
delay grows more slowly, because patients offered the longest waits are
precisely those who balk, selectively trimming the high-$\tau$ tail from the
accepted sample. By $\lambda = 110$ the gap reaches roughly $1.1$~days;
offered delay crosses the balking threshold while accepted delay has not yet
done so, illustrating why accepted wait understates congestion in
high-balking regimes.

\mainfigure{no_balk_delay_thresholds.png}{%
  Mean offered (dashed) and accepted (solid) booking delay from $\lambda = 30$
  to $\lambda = 110$, plotted against the no-show threshold ($6$~days, green)
  and balking threshold ($9$~days, purple). Shading marks the no-balking
  ($\lambda \leq 50$, blue) and balking ($\lambda \geq 60$, purple) regimes.
  The two series coincide where no one balks and diverge thereafter, with the
  gap widening as balking removes the high-$\tau$ tail from the accepted
  sample.
}{fig:no-balk-delay}

\subsection{Offered Wait vs.\ Accepted Wait}

The right panel of Figure~\ref{fig:main-panel} shows that accepted wait
is consistently below offered wait across all demand levels, with the gap
widening as balking rate rises; Figure~\ref{fig:no-balk-delay} shows the
same split in absolute terms as a function of~$\lambda$.

The implication is that accepted wait understates congestion in
high-balking regimes. Offered wait is the better denominator for any
congestion diagnostic---but it still excludes no-offer patients, who
receive no delay information at all.

\subsection{No-Offer as a Distinct Access Failure}

Mean offered wait is computed among patients who received a slot offer.
When the horizon saturates ($\lambda \gtrsim 110$ in this formulation),
a growing share of arrivals receive no offer and are entirely absent from
the offered-wait calculation. Figure~\ref{fig:no-balk-delay} covers
$\lambda \leq 110$, the range immediately preceding this saturation
regime; at $\lambda = 110$ offered delay has just crossed the balking
threshold and no-offer patients are beginning to accumulate. A useful
diagnostic is to report no-offer share alongside offered wait. Together
they describe both the level and the completeness of delay information.

%% ============================================================
\section{Sensitivity Screen: Main Patterns}
\label{sec:sensitivity}
%% ============================================================

The one-at-a-time sensitivity runs and the regression screen are described
in Appendix~\ref{app:sens-design} and~\ref{app:regression}. The patterns
most relevant for interpretation are summarized here.

\paragraph{Demand pressure.}
Total arrival rate has the largest absolute standardized regression
coefficient for overall served rate ($\beta = -0.774$) and is the
dominant driver of offered wait ($\beta = +0.576$). This is mechanically
consistent: more arrivals compete for a fixed number of slots, pushing
offers farther out and reducing access.

\paragraph{No-show and cancellation.}
Average no-show threshold has the largest coefficient for utilization
($\beta = +0.515$): a higher threshold delays the onset of high no-show
probability, keeping more booked slots as completed visits. Cancellation
has a positive coefficient for served rate ($\beta = +0.116$) and a
negative one for offered wait ($\beta = -0.442$). The positive sign on
served rate is counterintuitive at first: the mechanism is that
cancellations free future capacity, which can absorb more patients and
slightly improve the served fraction---a rebooking-chain effect under FCFS.

The mechanism is confirmed by the engine order and the arrival outcome
decomposition (Figure~\ref{fig:cancel-slice}). Cancellations fire at the
start of each day (step~1 of the daily loop), freeing future slots before
new arrivals are generated. New arrivals then book in FCFS order (step~3--5
of the same day), so every freed slot is immediately available for
rebooking---no slot goes empty as a result of cancellation alone. The net
effect on overall served rate is positive because: (i)~the replacement
arrival often books a nearer slot, reducing its no-show exposure; and
(ii)~at baseline load there is always a surplus of unserved arrivals to fill
freed capacity.

The class-level decomposition in Figure~\ref{fig:cancel-slice} makes the
asymmetry explicit. As one class's cancellation probability rises, that
class's own served rate falls sharply---its patients cancel themselves out.
The other class's served rate rises by a comparable amount, capturing the
freed slots. The net effect on overall served rate is small but consistently
positive across the full cancellation range. The effect would not hold at
very low demand, where freed slots might go unfilled; the positive
coefficient reflects the high-load regime that dominates the regression
sample.

\mainfigure{cancellation_probability_slice_access.png}{%
  Served rate as a function of cancellation probability (one class varies,
  the other held at baseline $\phi = 0.10$). The canceling class loses its
  own served rate; the other class gains by filling freed slots. Overall
  served rate (circles) rises modestly, confirming the rebooking-chain
  mechanism. Dashed vertical line marks the baseline.
}{fig:cancel-slice}

\paragraph{Balking and no-show step.}
The reparametrization into level and difference reveals two distinct
effects. The combined step level---the average of the balking and no-show
high-delay probabilities---carries a negative coefficient for utilization
($\beta = -0.288$), served rate ($\beta = -0.106$), and offered wait
($\beta = -0.181$). Higher behavioral intensity in either direction reduces
slot fill and access; the wait reduction is a selection effect (high-step
patients exit the queue), not a congestion improvement.

\paragraph{Balk vs.\ no-show dominance.}
The difference (balking step minus no-show step) is significant for
utilization ($\beta = +0.400$) and served rate ($\beta = +0.186$). When
balking exceeds no-show, patients who cannot obtain near-term slots
preemptively reject rather than accepting a distant booking and then
missing it. The net effect is that fewer booked slots go to waste: the
replacement arrival who takes a freed near-term slot tends to have lower
no-show exposure, while a no-show-dominant regime wastes already-booked
capacity with no equivalent rebooking benefit. The sign alone is not
significant in any target; the effect is driven by the magnitude of the
dominance, not a threshold.

Mediation analysis confirms that this path operates almost entirely through
the no-show rate: the coefficient on \texttt{balk\_minus\_noshow\_step} drops
from $0.284$ to $0.088$ once no-show rate is added as a covariate, with a
residual direct effect remaining. Figure~\ref{fig:balk-noshow-heatmap}
shows this structure directly. The heatmap gradient for utilization runs
almost entirely along the no-show step axis: rows at the top (high no-show
step) are uniformly dark regardless of balking step, while the steepest gains
from balking step appear only in the lower half of the heatmap where no-show
step is modest. The top projection (fix no-show step~$= 0.3$, vary balking
step) confirms the near-flat response of utilization to balking step alone.
The bottom projection (fix balking step~$= 0.5$, vary no-show step) shows
the steep monotone decline that no-show step imposes. The diagonal
(\texttt{balk\_step}~$=$~\texttt{noshow\_step}) marks the boundary between
the two dominance regimes; the baseline parameter point (balking
step~$= 0.50$, no-show step~$= 0.30$) sits below it, in balk-dominant
territory. The deep-dive analysis (Section~\ref{sec:balk-deepdive}) confirms
that accepted wait can fall while served rate also falls, consistent with the
level effect.

\mainfigure{balk_noshow_step_utilization.png}{%
  Left: average utilization as a function of balking step (x-axis) and
  no-show step (y-axis), both classes symmetric.  Dashed horizontal and
  dotted vertical lines mark the slice positions for the right panels; the
  diagonal marks \texttt{balk\_step}~$=$~\texttt{noshow\_step}.  Top right:
  utilization and no-show rate as balking step varies with no-show step fixed
  at~$0.3$---nearly flat, confirming that the balking step has little direct
  effect on utilization once no-show step is held constant.  Bottom right:
  the same metrics as no-show step varies with balking step fixed at~$0.5$---a
  steep monotone decline driven by wasted booked capacity.
}{fig:balk-noshow-heatmap}

\paragraph{Threshold ordering.}
Reparametrizing the two thresholds into level and difference reveals a
structurally distinct effect from the step analysis. The average threshold
level carries a positive coefficient for utilization ($\beta = +0.296$),
served rate ($\beta = +0.139$), and offered wait ($\beta = +0.148$):
pushing both thresholds out in days reduces the share of accepted patients
exposed to behavioral risk zones. The difference (balking threshold minus
no-show threshold) carries a significant \emph{negative} coefficient for
utilization ($\beta = -0.259$) and served rate ($\beta = -0.158$), and a
positive one for offered wait ($\beta = +0.173$).

The mechanism is structural. When balking threshold exceeds no-show
threshold (as at baseline: $\theta^{\text{balk}} = 9 > \theta^{\text{no-show}}
= 6$), a \emph{toxic zone} opens for $\tau \in (\theta^{\text{no-show}},\,
\theta^{\text{balk}}]$: patients in this delay range accept offers (no
balking risk yet) but carry positive no-show probability. These slots are
booked and then wasted at rate~$\xi^{\text{high}}$. When the thresholds
invert ($\theta^{\text{no-show}} > \theta^{\text{balk}}$), the toxic zone
vanishes: patients reach the balking threshold before the no-show threshold,
so those who would have been in the high-no-show zone either reject early
(balk) or, if they accept, have not yet entered the no-show region.
Crucially, the sign of the difference is itself significant for utilization
($\beta = -0.198$, $p = 0.003$), indicating a qualitative regime change at
the crossover---not merely a continuous slope. This contrasts with the step
dominance effect, where the sign was not significant.

Figure~\ref{fig:balk-noshow-threshold-heatmap} shows the heatmap and
projections. The gradient is strong along both axes but asymmetric about
the diagonal: above the diagonal (no-show threshold~$>$~balking threshold)
utilization is higher for a given average threshold level. The top
projection (fix no-show threshold at~$6$~d, vary balking threshold) is
nearly flat---once the no-show threshold is held fixed, moving the balking
threshold around it has little additional effect. The bottom projection
(fix balking threshold at~$9$~d, vary no-show threshold) shows a steep
rise as the no-show threshold increases past the balking threshold,
confirming the regime-change interpretation.

\mainfigure{balk_noshow_threshold_utilization.png}{%
  Left: average utilization as a function of balking threshold (x-axis) and
  no-show threshold (y-axis), both thresholds symmetric across classes.
  The filled circle marks the baseline ($\theta^{\text{balk}} = 9$~d,
  $\theta^{\text{no-show}} = 6$~d), which lies below the diagonal in
  the toxic-zone regime. Above the diagonal the no-show threshold exceeds
  the balking threshold, eliminating the accept-then-miss zone and raising
  utilization. Top right: fix no-show threshold at~$6$~d---nearly flat as
  balking threshold varies, confirming that the toxic zone width (not the
  balking threshold level alone) drives the effect. Bottom right: fix
  balking threshold at~$9$~d---utilization rises steeply as no-show
  threshold increases past the balking threshold (at the dashed vertical
  line), then flattens once the toxic zone is closed.
}{fig:balk-noshow-threshold-heatmap}

\paragraph{Class gap drivers.}
In the randomized screen, the largest class-gap drivers are cancellation
probability gap ($|\beta| = 0.452$), balking threshold gap ($|\beta| =
0.429$), and balking step gap ($|\beta| = 0.349$). These findings are
plausible: in a pooled FCFS system, class disparities arise because one
class accepts more long-delay offers or is more likely to cancel after
booking.

\mainfigure{regression_significant_standardized_coefficients.png}{%
Significant standardized regression coefficients from 240 randomized
parameter settings. Features and targets are standardized. The
dominance of total arrival rate for served rate and offered wait is
consistent with the structural decoupling of utilization and access.
Error bars are 95\% HC3 robust confidence intervals. Full table in
Appendix~\ref{app:regression}.}{fig:regression}

%% ============================================================
\section{Balking Deep Dive: Selection vs.\ Access}
\label{sec:balk-deepdive}
%% ============================================================

This section isolates Class~1 balking while Class~2 stays at baseline, to
test whether lower accepted wait can accompany lower access.

\mainfigure{../../../../../outputs/class1_balking/figures/percent_serviced_by_class.png}{%
As Class~1 high-balking probability increases from 0 to~1, Class~1
served rate falls while Class~2 served rate rises. The congestion
relief that Class~2 gains is mechanically consistent with fewer
Class~1 patients booking long-delay slots.}{fig:balk-served}

\mainfigure{../../../../../outputs/class1_balking/figures/mean_accepted_delay_by_class.png}{%
Class~1 accepted wait falls as Class~1 high-balking probability
increases. This is a selection artifact: long-delay Class~1 patients
leave the accepted pool. Accepted wait should not be read as a
delay-improvement signal in this context.}{fig:balk-accepted}

The two figures together demonstrate a pattern the formulation produces
cleanly: accepted wait and served rate move in opposite directions
when balking changes. Any policy that raises balking step would show
apparently improved wait times while actually reducing the share of
Class~1 patients who complete service.

%% ============================================================
\section{Class-Gap Interpretation Under Pooled FCFS}
\label{sec:classgap}
%% ============================================================

Arrivals from both classes are pooled and randomly permuted each day before
FCFS assignment. The engine does not distinguish between classes when
offering slots. Therefore:

\begin{itemize}[leftmargin=1.5em, itemsep=3pt]
\item When class parameters are symmetric, the expected class served-rate
gap is zero. Observed deviations are simulation noise.
\item Meaningful class gaps arise only from asymmetric behavioral parameters
(different balking threshold, balking step, cancellation probability, or
no-show behavior) or from different arrival rates. These create systematic
differences in how often each class successfully books, cancels, and
attends.
\item A positive class~1 balking threshold relative to class~2 means
class~1 tolerates longer delays before rejecting offers. Under FCFS, a
more tolerant class captures more of the later-horizon slots and therefore
achieves a higher served rate.
\end{itemize}

\mainfigure{metric_class_gap_drivers.png}{%
When Class~1 behavioral parameters are varied while Class~2 stays at
baseline, a gap opens. Higher Class~1 cancellation, higher balking step,
or higher no-show step widen the gap against Class~1. Higher Class~1
balking threshold helps Class~1 by allowing it to accept longer-delay
slots. With symmetric parameters, both class lines coincide.
}{fig:class-gap}

The regression screen confirms this ranking: among class-gap drivers,
cancellation probability gap dominates ($|\beta| = 0.452$), followed by
balking threshold gap ($|\beta| = 0.429$) and balking step gap
($|\beta| = 0.349$). Class~1 arrival share enters significantly
($\beta = -0.120$): a larger Class~1 share tends to mildly reduce the
Class~1 served-rate advantage, consistent with more Class~1 arrivals
competing for the same slots under FCFS.

%% ============================================================
\section{What Seems Mechanically Robust vs.\ What Needs Validation}
\label{sec:checks}
%% ============================================================

\subsection{Mechanically Robust Under the Current Formulation}

\begin{itemize}[leftmargin=1.5em, itemsep=3pt]
\item \textbf{Utilization and served rate have different denominators.}
Their decoupling is a direct consequence of the metric definitions. It
is not a fragile numerical result.

\item \textbf{High utilization can coexist with low served rate.}
When demand exceeds supply, served rate is bounded by the supply-to-demand
ratio. This bound holds regardless of behavioral parameters.

\item \textbf{Offered wait is less selective than accepted wait.}
Offered delay is recorded before balking, so offered wait includes patients
who ultimately leave. Accepted wait excludes them. The direction of the
difference is structural.

\item \textbf{No-offer is a separate access failure not captured by
offered wait.} No-offer patients receive no delay information and are
excluded from all delay metrics. Above $\lambda \approx 110$, they become
a substantial fraction of arrivals.

\item \textbf{Class labels are exchangeable under symmetric parameters.}
FCFS pooling and random permutation mean that class identity has no
systematic effect when classes are identical. The near-zero baseline class
gap (0.001) is consistent with this.

\item \textbf{The cancellation--delay feedback is directionally correct.}
Longer average booking delay creates more cancellation-check opportunities,
which raises cumulative cancellation probability. The high observed
cancellation fraction at long delays is mechanically consistent with this.

\item \textbf{No-offer is zero at baseline; access failures are behavioral.}
The booking horizon does not saturate at $\lambda = 100$ because
cancellations continuously free future capacity. Balking and cancellation,
not horizon saturation, are the primary failure modes.
\end{itemize}

\subsection{Needs Validation or Further Code Checks}

\vspace{6pt}
\begin{checkenv}
\begin{enumerate}[leftmargin=1.5em, itemsep=4pt]
\item \textbf{Confirm cancellation applies once per residual day.}
The documentation says cancellation is checked each day for all future
appointments. If the engine instead applies cancellation once at booking,
the high cancellation fraction is unexplained and a code audit is needed.
Check \texttt{simulation/engine.py} cancellation loop.

\item \textbf{Clarify no-offer propagation within a day.}
The documentation states that when the horizon is full for one patient,
all remaining arrivals that day are counted as \texttt{no\_offer}. This
means a single saturation event affects the full day remainder, which
could create day-order artifacts. Verify whether the no-offer share is
sensitive to the random permutation seed.

\item \textbf{Assess whether no-show slots should remain unbooked.}
No-show slots are not rebooked in the current formulation. In many
real clinics, same-day no-shows are filled with walk-ins or advance
overbooking. If this assumption is changed, both utilization and served
rate would shift; the magnitude depends on the no-show fraction and
same-day availability.

\item \textbf{Assess whether same-day cancellations should be excluded.}
Currently, no same-day cancellations are allowed. This keeps service-day
slots deterministic once the service day is reached. Relaxing this would
add a slot-recapture mechanism that could affect both utilization and
served rate.

\item \textbf{Confirm stability of class-gap estimates.}
The class-gap metric has the lowest out-of-sample $R^2$ (0.708) in the
regression screen and is near zero (0.001) at symmetric baseline. With
only 2~seeds per setting in the regression screen, sampling variance in
the gap could be material. At least 10--20 seeds per grid point are
advisable before drawing strong conclusions about class equity.

\item \textbf{Evaluate realism of baseline demand and capacity levels.}
100 arrivals per day against 32 slots per day is an intentional
stress-test, not a calibrated clinic. Before drawing any operational
conclusions, behavioral parameters (balking thresholds, no-show
probabilities, cancellation rates) should be estimated from clinic data
and the demand-to-capacity ratio brought into a plausible range.
\end{enumerate}
\end{checkenv}

%% ============================================================
\section{Summary for Discussion}
\label{sec:summary}
%% ============================================================

\begin{enumerate}[leftmargin=1.5em, itemsep=4pt]
\item Under the current formulation, utilization and served rate are
structurally decoupled. At baseline, this produces a 3:1 discrepancy
(0.839 vs.\ 0.269). Using utilization as a proxy for access would be
misleading at this demand level.

\item The dominant loss channels at baseline are balking (35.6\%) and
cancellation (32.3\%). No-offer is zero. Above $\lambda \approx 110$,
horizon saturation becomes an additional channel.

\item The high cancellation fraction is consistent with the long average
booking delay and the daily-cancellation rule. Whether this is realistic
depends on how real patients behave when booked many days out.

\item Offered wait is preferable to accepted wait as a delay measure,
because it includes patients who rejected the offer. But it still excludes
no-offer patients, who are the most severely access-constrained.

\item Class disparities in the simulation arise from asymmetric behavioral
parameters, not from FCFS favoritism. Symmetric classes produce
near-zero gaps.

\item The regression screen and one-at-a-time sensitivity runs are
internally consistent: demand pressure drives served rate and offered
wait; no-show and cancellation behavior drive utilization; behavioral
parameter gaps drive class equity.

\item None of these findings should be interpreted as causal estimates
about real appointment systems. They describe how this particular model
formulation responds to parameter variation.
\end{enumerate}

\bigskip
\noindent\textbf{Questions for discussion:}
\begin{enumerate}[leftmargin=1.5em, itemsep=4pt]
\item Is the baseline demand/capacity ratio intended as a stress scenario,
or should it be recalibrated toward a more plausible clinic setting?
\item Should cancellation be modeled as a daily hazard (once per residual
day) or a one-time event at booking? The cumulative cancellation fraction
depends critically on this.
\item Should no-offer patients be reported as a separate access failure
in all summary tables, not only when the no-offer rate is non-negligible?
\item Should no-show probability depend on the original offered delay
$\tau$ (current formulation) or on the remaining wait at service time?
\item Should same-day cancellations be allowed? What would the rebooking
mechanism look like?
\item Should alternative scheduling rules (e.g., priority-based,
overbooking) be compared against FCFS within this formulation?
\end{enumerate}

%% ============================================================
\appendix

\section{Exact Metric Definitions and Notation}
\label{app:defs}
%% ============================================================

Let $i$ be patient class, $D$ measured service days, $S$ slots per day,
$A_i$ arrivals, $B_i$ accepted bookings, $K_i$ balked patients, $C_i$
cancellations, $H_i$ no-shows, $V_i$ completed visits, and
$O_i = B_i + K_i$ offered patients.

\begin{align*}
\texttt{percent\_serviced}_i &= V_i / A_i, \\
\texttt{overall\_percent\_serviced} &= \textstyle\sum_i V_i \;/\; \textstyle\sum_i A_i, \\
\texttt{mean\_offered\_booking\_delay} &= \textstyle\sum_i \tau^{\text{off}}_i \;/\; \textstyle\sum_i O_i, \\
\texttt{mean\_accepted\_booking\_delay} &= \textstyle\sum_i \tau^{\text{acc}}_i \;/\; \textstyle\sum_i B_i, \\
\texttt{average\_utilization} &= \tfrac{1}{D}\textstyle\sum_d V_d / S, \\
\texttt{balking\_rate} &= \textstyle\sum_i K_i \;/\; \textstyle\sum_i O_i, \\
\Delta_{\text{access}} &= \texttt{percent\_serviced}_1 - \texttt{percent\_serviced}_2.
\end{align*}

\noindent\textbf{Note on terminology.} In the main text,
\texttt{percent\_serviced} (the simulation's internal variable name) is
reported as \emph{overall served rate}. \texttt{mean\_offered\_booking\_delay}
is reported as \emph{mean offered wait}; \texttt{mean\_accepted\_booking\_delay}
as \emph{mean accepted wait}.

\medskip
\noindent\textbf{Behavioral step functions.}
\[
  b_i(\tau) = \begin{cases}
    b_i^{\text{low}} & \tau < \theta_i^{\text{balk}}, \\
    b_i^{\text{high}} & \tau \geq \theta_i^{\text{balk}}.
  \end{cases}
  \qquad
  \xi_i(\tau) = \begin{cases}
    \xi_i^{\text{low}} & \tau < \theta_i^{\text{no-show}}, \\
    \xi_i^{\text{high}} & \tau \geq \theta_i^{\text{no-show}}.
  \end{cases}
\]

Balking step $= b_i^{\text{high}} - b_i^{\text{low}}$; no-show step
$= \xi_i^{\text{high}} - \xi_i^{\text{low}}$.

\medskip
\noindent\textbf{Measurement semantics.} Class metrics cover patients
arriving in the measurement window. If \texttt{cooldown\_days}~$<$~%
\texttt{horizon\_days}~$-$~1, some late bookings are unresolved at
simulation end and counted as \texttt{unresolved\_booked} (0.06\% of
arrivals at baseline). \texttt{total\_value} is service-day based and
is not cohort-aligned with class metrics.

%% ============================================================
\section{Simulation Design and Sensitivity Grid}
\label{app:sens-design}
%% ============================================================

All sensitivity grids are centered on \texttt{configs/baseline.yaml}.
One-axis slices vary Class~1 parameters while Class~2 stays at baseline.
Two-class heatmaps vary both classes independently.

\begin{table}[H]
\centering
\renewcommand{\arraystretch}{1.15}
\begin{tabular}{p{0.30\textwidth}rrr}
\toprule
Parameter & Baseline & Range & Grid size \\
\midrule
Arrivals (per class), 1D slices & 50 & 25--75 & continuous \\
Total arrivals, capacity stress & 100 & 30--170 & continuous \\
Balking step & 0.50 & 0.00--1.00 & 21 \\
Balking threshold (days) & 9 & 0--13 & 14 \\
No-show step & 0.30 & 0.00--1.00 & 21 \\
No-show threshold (days) & 6 & 0--13 & 14 \\
Cancellation probability & 0.10 & 0.00--0.30 & 21 \\
\bottomrule
\end{tabular}
\caption{Sensitivity parameter ranges. Each grid point uses seed 2027;
multi-seed comparisons average seeds 2027, 2028, 2029.}
\label{tab:sens-grid}
\end{table}

\textbf{Regression screen.}
240 randomized parameter settings (seed 9137); two seeds (4101, 4102)
averaged per setting; 12 features (total arrival rate, Class~1 share,
average level and Class~1$-$Class~2 gap for each behavioral parameter);
OLS with HC3 robust standard errors; only $p < 0.05$ coefficients shown.

\textbf{Balking deep dives.}
Class~1 high-balking-probability sweep over $\{0.0, 0.1, \ldots, 0.9\}$
(100 seeds per point; Class~2 fixed at baseline). Class~1 threshold sweep
over $\{3, 4, \ldots, 13\}$ (100 seeds per point; Class~2 fixed at
baseline).

%% ============================================================
\section{Supplementary Sensitivity Plots}
\label{app:sens}
%% ============================================================

\subsection{Driver Plots}

\mainfigure{metric_utilization_drivers.png}{%
Utilization across Class~1 arrival and behavioral parameter ranges.
Utilization stays near 0.839 across most of the demand range, while
class-level served rates diverge.}{fig:util-drivers}

\mainfigure{metric_access_drivers.png}{%
Served rate falls as arrival rates rise. The Class~1 and Class~2 lines
separate when Class~1 parameters diverge from symmetric baseline.
}{fig:access-drivers}

\mainfigure{metric_wait_drivers.png}{%
Offered and accepted wait across Class~1 parameter ranges. Accepted
wait can fall while offered wait stays elevated, illustrating selection
bias in the accepted-wait metric.}{fig:wait-drivers}

\mainfigure{metric_balking_rate_drivers.png}{%
Balking rate is controlled primarily by balking step and threshold.
Rising balking rate indicates congestion but not improved service.
}{fig:balk-drivers}

\subsection{No-Show and Cancellation Heatmaps}

\mainfigure{no_show_step_interaction_heatmap.png}{%
No-show step heatmap. Utilization falls when both classes carry high
no-show risk. Class gap moves against the higher no-show class.
}{fig:noshow-step}

\mainfigure{no_show_threshold_interaction_heatmap.png}{%
No-show threshold heatmap. Higher threshold delays high no-show onset,
improving utilization and access for the affected class.
}{fig:noshow-threshold}

\mainfigure{cancellation_probability_interaction_heatmap.png}{%
Cancellation heatmap. Higher cancellation for one class reduces that
class's access; aggregate effects are more complex (see
Section~\ref{sec:sensitivity}).}{fig:cancel}

\subsection{Balking Heatmaps}

\mainfigure{balking_step_interaction_heatmap.png}{%
Balking step heatmap. Served rate falls and class gap widens as one
class's balking step rises above the other.}{fig:balk-step}

\mainfigure{balking_threshold_interaction_heatmap.png}{%
Balking threshold heatmap. The more tolerant class (higher threshold)
captures more served slots under FCFS.}{fig:balk-threshold}

\mainfigure{arrival_mix_interaction_heatmap.png}{%
Arrival mix heatmap. Left: overall served rate. Right: mean offered wait.}{fig:arrival-mix}

\subsection{Capacity Stress}

\mainfigure{fcfs_capacity_stress_curves.png}{%
Capacity stress diagnostic curves. Served rate and offered wait as
demand rises from $\lambda = 30$ to~170; utilization remains nearly
flat across the full range.}{fig:stress}

\subsection{Balking Deep Dive: Additional Plots}

\mainfigure{../../../../../outputs/class1_balking/figures/mean_offered_delay_by_class.png}{%
Offered wait by class as Class~1 high-balking probability rises.
When enough Class~1 patients reject long-delay offers, overall horizon
congestion falls and both classes receive shorter offered delays.
}{fig:balk-offered}

\mainfigure{../../../../../outputs/class1_balking_threshold/figures/percent_serviced_by_class.png}{%
Class~1 balking threshold sweep: higher threshold increases Class~1
access by tolerating longer delays. The gain comes at mild cost to
Class~2 when Class~1 tolerance is very high.}{fig:thresh-served}

\mainfigure{../../../../../outputs/class1_balking_threshold/figures/service_gap.png}{%
Class~1 minus Class~2 served-rate gap as Class~1 threshold varies.
}{fig:thresh-gap}

%% ============================================================
\section{Regression Screen}
\label{app:regression}
%% ============================================================

\begin{table}[H]
\centering
\small
\renewcommand{\arraystretch}{1.08}
\begin{tabular}{p{0.17\textwidth}p{0.27\textwidth}rp{0.20\textwidth}r}
\toprule
Target & Feature & Std.\ coef. & 95\% HC3 CI & HC3 $p$ \\
\midrule
Utilization & Avg threshold level (balk $+$ no-show) & $+0.296$ & $[0.219,\;0.373]$ & $<0.001$ \\
Utilization & Balk step $-$ no-show step & $+0.365$ & $[0.227,\;0.504]$ & $<0.001$ \\
Utilization & Avg cancellation prob. & $+0.311$ & $[0.240,\;0.384]$ & $<0.001$ \\
Utilization & Avg step level (balk $+$ no-show) & $-0.289$ & $[-0.371,\;-0.207]$ & $<0.001$ \\
Utilization & Balk threshold $-$ no-show threshold & $-0.259$ & $[-0.387,\;-0.131]$ & $<0.001$ \\
Utilization & Sign(balk threshold $-$ no-show threshold) & $-0.198$ & $[-0.330,\;-0.067]$ & $0.003$ \\
Utilization & Total arrival rate & $-0.104$ & $[-0.182,\;-0.027]$ & $0.008$ \\
\midrule
Overall served rate & Total arrival rate & $-0.773$ & $[-0.854,\;-0.692]$ & $<0.001$ \\
Overall served rate & Avg threshold level (balk $+$ no-show) & $+0.139$ & $[0.084,\;0.195]$ & $<0.001$ \\
Overall served rate & Balk step $-$ no-show step & $+0.174$ & $[0.077,\;0.271]$ & $<0.001$ \\
Overall served rate & Avg step level (balk $+$ no-show) & $-0.107$ & $[-0.168,\;-0.045]$ & $0.001$ \\
Overall served rate & Balk threshold $-$ no-show threshold & $-0.158$ & $[-0.264,\;-0.053]$ & $0.003$ \\
Overall served rate & Avg cancellation prob. & $+0.114$ & $[0.054,\;0.174]$ & $<0.001$ \\
\midrule
Mean offered wait & Total arrival rate & $+0.575$ & $[0.496,\;0.653]$ & $<0.001$ \\
Mean offered wait & Avg cancellation prob. & $-0.440$ & $[-0.504,\;-0.375]$ & $<0.001$ \\
Mean offered wait & Balk threshold $-$ no-show threshold & $+0.173$ & $[0.049,\;0.296]$ & $0.006$ \\
Mean offered wait & Avg threshold level (balk $+$ no-show) & $+0.148$ & $[0.069,\;0.228]$ & $<0.001$ \\
Mean offered wait & Avg step level (balk $+$ no-show) & $-0.181$ & $[-0.271,\;-0.091]$ & $<0.001$ \\
\midrule
Class gap & Cancellation prob.\ gap & $-0.452$ & $[-0.558,\;-0.346]$ & $<0.001$ \\
Class gap & Balking threshold gap & $+0.429$ & $[0.331,\;0.528]$ & $<0.001$ \\
Class gap & Balking step gap & $-0.352$ & $[-0.436,\;-0.268]$ & $<0.001$ \\
Class gap & No-show threshold gap & $+0.266$ & $[0.175,\;0.357]$ & $<0.001$ \\
Class gap & No-show step gap & $-0.213$ & $[-0.289,\;-0.136]$ & $<0.001$ \\
Class gap & Class~1 arrival share & $-0.119$ & $[-0.209,\;-0.029]$ & $0.009$ \\
\bottomrule
\end{tabular}
\caption{Significant standardized regression coefficients ($p<0.05$, HC3 robust SEs).
Gap = Class~1 minus Class~2 value. 240 randomized settings, 2 seeds averaged per
setting. Only features passing $p < 0.05$ are shown.}
\label{tab:regression}
\end{table}

\mainfigure{regression_significant_standardized_coefficients.png}{%
Significant standardized regression coefficients with 95\% HC3 confidence
intervals. Total arrival rate dominates served rate and offered wait;
no-show behavior dominates utilization; behavioral parameter gaps drive
the class served-rate gap.}{fig:app-regression}

%% ============================================================
%% Changelog
%% ============================================================
% Base file: metric_analysis_research_note.tex (v1)
% Sections rewritten:
%   - §2 Baseline (added formal parameter table)
%   - §10 Summary (renamed "Summary for Discussion"; added discussion questions)
% Sections added:
%   - §4 Cross-Metric Diagnostics (from merged)
%   - App B Sensitivity Design (from merged)
%   - App D Regression Screen (from merged)
% Figures in main body:
%   - fcfs_arrival_outcome_decomposition.png (§3)
%   - research_note_main_panel.png (§4)
%   - regression_significant_standardized_coefficients.png (§6)
%   - class1_balking deep-dive figures (§7)
%   - metric_class_gap_drivers.png (§8)
% Figures moved to appendix:
%   - metric_utilization_drivers.png, metric_access_drivers.png,
%     metric_wait_drivers.png, metric_balking_rate_drivers.png (App C)
%   - all heatmaps (App C)
%   - fcfs_capacity_stress_curves.png (App C)
%   - additional deep-dive threshold figures (App C)
% Figures removed:
%   - scenario_metric_comparison.png (Scenario 2 reference)
% Scenario 2: Removed entirely.
%   Reason: configs/scenario_2.yaml changes 6 parameters simultaneously
%   (horizon, lambda_2, balk_threshold, balk_high, no_show_threshold,
%   no_show_high). Not a one-parameter comparison. util=1.000 in merged
%   table is suspicious. Per hard constraint #2, removed.
% Unresolved validation checks (carried as checkenv callouts):
%   1. Cancellation: once per residual day vs. once per booking
%   2. No-offer propagation creates daily-order artifacts
%   3. Positive cancellation coefficient on served rate
%   4. Class-gap stability with more seeds

\end{document}
